\chapter{NHẬP MÔN: TỪ PIXEL ĐẾN NHẬN THỨC}
\label{chap:nhap_mon}

Mỗi khoảnh khắc mắt chúng ta mở ra là một bản giao hưởng thị giác được trình diễn. Trong một cái chớp mắt, não bộ con người có thể nhận diện khuôn mặt thân quen giữa đám đông, cảm nhận chiều sâu của một con đường, phân biệt hàng ngàn màu sắc, và thậm chí đọc được cảm xúc tinh tế trên nét mặt của người đối diện. Khả năng "nhìn" và "hiểu" này là một trong những kỳ công vĩ đại nhất của tạo hóa, một quá trình diễn ra tức thời, tự nhiên đến mức chúng ta hiếm khi dừng lại để suy ngẫm về sự phức tạp phi thường đằng sau nó.

Nhưng đối với một cỗ máy, thế giới thị giác lại khởi đầu hoàn toàn khác. Nơi con người nhìn thấy một bông hoa, máy tính chỉ thấy một ma trận số khổng lồ, vô hồn. Mỗi con số đại diện cho cường độ sáng của một điểm ảnh duy nhất—một pixel. Làm thế nào để từ một biển số liệu khô khan đó, máy tính có thể đi đến kết luận rằng "đây là một bông hồng đỏ"? Làm thế nào nó có thể phân biệt được bông hồng đó với chiếc lá xanh bên cạnh, hay nhận ra rằng nhiều bông hồng hợp lại tạo thành một khu vườn?

Chương đầu tiên này chính là câu chuyện về hành trình kỳ diệu đó: hành trình \textbf{từ pixel đến nhận thức}. Đây là cuộc phiêu lưu trí tuệ nhằm giải mã ngôn ngữ của ánh sáng và màu sắc, dạy cho máy móc cách diễn giải thế giới ba chiều từ những hình ảnh hai chiều phẳng lặng. Chúng ta sẽ bắt đầu từ những đơn vị cơ bản nhất, những "nguyên tử" của thế giới ảnh số, và từng bước xây dựng nên các khái niệm phức tạp hơn.

Điểm khởi đầu của chúng ta là \textbf{pixel}, một chấm màu đơn độc, không mang ý nghĩa. Nhưng khi hàng triệu pixel kết hợp lại, chúng tạo nên những cấu trúc, những hoa văn. Chúng ta sẽ học cách nhóm các pixel lại thành các đường biên, các vùng màu sắc, các hình khối. Dần dần, từ những cấu trúc đơn giản này, chúng ta sẽ trích xuất ra những đặc trưng có ý nghĩa hơn, những "từ vựng" trong ngôn ngữ thị giác của máy móc. Và cuối cùng, khi kết hợp những từ vựng đó lại, máy tính sẽ bắt đầu "hiểu" được ngữ nghĩa của cảnh vật—nhận ra sự hiện diện của con người, xe cộ, cây cối. Đó chính là đỉnh cao của hành trình: \textbf{nhận thức}.

Nhận thức không chỉ là nhận dạng. Đó là sự thấu hiểu. Đó là khả năng của một chiếc xe tự lái không chỉ "thấy" một vật thể màu đỏ phía trước, mà còn "hiểu" rằng đó là đèn tín hiệu giao thông và cần phải dừng lại. Đó là khả năng của một hệ thống y tế không chỉ "nhìn" thấy các đốm mờ trên ảnh X-quang, mà còn "hiểu" rằng đó có thể là dấu hiệu sớm của một khối u.

Hành trình này không còn là khoa học viễn tưởng. Nó đang diễn ra xung quanh chúng ta, thay đổi bộ mặt của xã hội hiện đại. Từ những chiếc xe tự lái đang định hình lại ngành giao thông, các trợ lý AI có thể chẩn đoán bệnh với độ chính kim xác đáng kinh ngạc, cho đến các ứng dụng thực tế tăng cường biến thế giới thực thành một không gian tương tác sống động—tất cả đều bắt nguồn từ nguyên lý cốt lõi: biến pixel thành nhận thức.

Trong chương này, chúng ta sẽ xây dựng nền móng cho toàn bộ cuốn sách. Chúng ta sẽ:
\begin{itemize}
	\item Định nghĩa một cách chính xác \textbf{Thị giác Máy tính là gì} và khám phá phạm vi rộng lớn của nó.
	\item Ôn lại \textbf{ngôn ngữ toán học} cần thiết—đại số tuyến tính, giải tích, xác suất—những công cụ không thể thiếu để mô tả và thao tác với dữ liệu thị giác.
	\item Tìm hiểu cách \textbf{ảnh và video kỹ thuật số được biểu diễn} bên trong máy tính.
	\item Khám phá những \textbf{cảm biến hình ảnh tiên tiến} vượt ra ngoài camera thông thường, mở ra những cách "nhìn" mới cho máy móc.
	\item Và quan trọng không kém, chúng ta sẽ đối mặt với những câu hỏi về \textbf{đạo đức, sự công bằng và quyền riêng tư} trong kỷ nguyên mà máy móc có thể "nhìn" ở khắp mọi nơi.
\end{itemize}

Đây là chương bản lề, trang bị cho bạn những khái niệm, công cụ và tư duy cần thiết để sẵn sàng chinh phục những kiến trúc phức tạp và các thuật toán tiên tiến nhất trong thế giới Thị giác Máy tính. Hãy cùng nhau bắt đầu cuộc phiêu lưu vào thế giới thị giác của máy móc, một thế giới được dệt nên từ hàng tỷ pixel.
